\documentclass[10pt]{article}
\usepackage[margin=20mm]{geometry}
\usepackage{amsmath,amssymb,mathtools,fontspec,unicode-math,xurl,hyperref}
\IfFontExistsTF{DejaVu Serif}{\setmainfont{DejaVu Serif}}{\setmainfont{Latin Modern Roman}}
\setmathfont{latinmodern-math.otf}
\setlength{\parindent}{0pt}\setlength{\parskip}{6pt}\sloppy
\title{Technical verification of relations on moduli spaces and the negative \(r\)-spin Witten conjecture}
\author{AI-assisted technical review}
\date{}
\begin{document}
\maketitle
Abstract: One Major issue is established in the acquired publisher version: the universal string-and-reduction criterion in Proposition 5.6 is false at \(r=2\). A normalized stable formal series satisfies every written hypothesis but fails a required weight-two constraint. The proof uses a weight-two annihilation relation that its preceding induction does not supply at this endpoint. This criterion is used in the paper’s route to a principal integrability conclusion. The counterexample does not refute the particular Brezin–Gross–Witten identification; the paper itself cites full \(r=2\) constraints that may supply a separate proof. Three additional checks support the stated rank and normalization calculations, the unmarked top-degree exception, and finite coefficients of the displayed \(\kappa\)-series. All conclusions have the explicit coverage limits below. Scientific adjudication remains pending.\par
\section*{Source identity and scope}
The reviewed work is Nitin Kumar Chidambaram, Elba Garcia-Failde, and Alessandro Giacchetto, Relations on \(\overline{\mathcal M}_{g,n}\) and the negative \(r\)-spin Witten conjecture, Inventiones mathematicae 241 (2025), 929–997, DOI 10.1007/s00222-025-01351-y [paper]. The acquired 69-page PDF is the publisher Version of Record. Throughout this report, PDF page \(p\) is printed page \(p+928\); both locators are given for the decisive issue. No later correction, author adoption, or equivalence to another version is asserted.\par
The principal domain is an integer \(r\ge2\), stable pairs \((g,n)\) with \(2g-2+n>0\), and positive primary labels \(1\le a_i\le r-1\) for the \(\Theta\) theory. The auxiliary \(\Upsilon\) construction also uses zero primary labels. Theorem E/5.5 states the full \(W\)-constraints, whereas the general \(r\)-KdV/BGW identification is Conjecture F/5.17; Theorem G/5.20 proves the stated \(r=2,3\) cases. The review does not attribute an all-\(r\) integrability theorem to the authors.\par
The existing same-conversation native helper `/root/check\_0090` supplied bounded source checks and independently written exact symbolic calculations. A second existing helper, `/root/check\_0089`, challenged the hypotheses, normal ordering, and downstream scope of the proposed counterexample analytically. This is an initial technical review, not independent acceptance. Author software was not executed.\par
\section*{Major: the \(r=2\) endpoint of Proposition 5.6}
\subsection*{Claim, operators, and exact source locations}
Equation reference: eq:stable, eq:criterion. Proposition 5.6, equations (5.16)–(5.19), appears on PDF pp. 52–53, printed pp. 980–981. It asserts that a formal function of the form \[
 Z=\exp\!\left(
 \sum_{\substack{g\ge0,\ n\ge1\\2g-2+n>0}}
 \frac{\hbar^{g-1}}{n!}
 \sum_{a_1,\ldots,a_n>0}
 \Phi_{g,n}[a_1,\ldots,a_n]\prod_{j=1}^n t_{a_j}
 \right),\] with symmetric coefficient arrays, must equal \(Z^{\Theta^r}\) if \[
 H_k^1Z=\nu_kZ\quad(k\ge1),\qquad
 H_{-r+2}^{r}Z=\mu Z\] for constants \(\nu_k,\mu\). There is no \(r\ge3\) restriction or KP/KdV hypothesis in this proposition. The surrounding text introduces the assertion before discussing integrability. A later application to a particular KP tau function does not supply an additional premise to this universal statement.\par
Equation reference: eq:Hone, eq:Htwo. For reproducibility, the relevant definitions are (5.2) on PDF p. 48/printed p. 976, the shift (5.6)–(5.7) on PDF p. 49/printed p. 977, and (5.8) on PDF p. 50/printed p. 978: \[J_\ell=
 \begin{cases}
 \hbar\,\partial_{t_\ell},&\ell>0,\\
 -\ell\,t_{-\ell},&\ell\le0,
 \end{cases}
 \qquad J_0=0,\qquad
 \widetilde J_\ell=J_\ell-\delta_{\ell,-r+1},\] \[\begin{aligned}
 H_k^1&=\widetilde J_{rk},\\
 H_k^2&=\frac12\sum_{p_1+p_2=rk}
 \bigl(r\,\mathbf1_{r\mid p_1,\ r\mid p_2}-1\bigr)
 :\widetilde J_{p_1}\widetilde J_{p_2}:
 -\frac{\hbar(r^2-1)}{24}\delta_{k,0}.
 \end{aligned}\] The superscript \(2\) denotes conformal weight, not an operator square. The colons place multiplication operators before positive-index derivatives.\par
\subsection*{A formal series meeting every stated hypothesis}
Equation reference: eq:Zzero. Set \(r=2\) and take the binomial expansion at the origin \[
 Z_0=(1-t_1)^{-1/8},\qquad
 \log Z_0=\sum_{n\ge1}\frac{t_1^n}{8n}.\] Define \[\Phi_{1,n}[1,\ldots,1]=\frac{(n-1)!}{8}\quad(n\ge1),
 \qquad \Phi_{g,n}[a_1,\ldots,a_n]=0\] in every other case. Substitution into [eq:stable] (\#eq:stable) gives precisely [eq:Zzero] (\#eq:Zzero): the factor \(1/n!\) leaves \(1/(8n)\), and \(\hbar^{g-1}=1\). The arrays are symmetric; all indices are positive. Every nonzero term has \(g=1\), \(n\ge1\), and \(2g-2+n=n>0\). There is no \(n=0\) term, and \(Z_0(0)=1\). Thus neither instability nor an arbitrary multiplicative normalization explains the example. Formal expansion at the origin is sufficient; no analytic branch continuation is used.\par
Equation reference: eq:Hzero. The series is independent of every even time. Equation [eq:Hone] (\#eq:Hone) therefore gives \[H_k^1Z_0=\hbar\,\partial_{t_{2k}}Z_0=0
 \qquad(k\ge1),\] so \(\nu_k=0\). At \(r=2\), the shift is \(\widetilde J_{-1}=t_1-1\). On functions independent of the even times, [eq:Htwo] (\#eq:Htwo) reduces at mode zero to \[
 H_0^2=\hbar\left(
 (1-t_1)\partial_{t_1}
 -\sum_{\substack{j\ge3\\j\ {\rm odd}}}j t_j\partial_{t_j}
 -\frac18\right).\] Since \(\partial_{t_1}Z_0=\tfrac18(1-t_1)^{-9/8}\), this gives \(H_0^2Z_0=0\). The other hypothesis of [eq:criterion] (\#eq:criterion) holds with the allowed constant \(\mu=0\).\par
Even the geometric dimension selection rule would not discard these coefficients: at \(r=2\), all \(a_i=1\), the class degree is \(D=2g-2+n\), so the total descendant degree must be \(g-1\). The nonzero terms in [eq:Zzero] (\#eq:Zzero) have genus one and all descendant indices zero. This observation does not assume that \(Z_0\) arises from a CohFT, which is not a written hypothesis of Proposition 5.6.\par
\subsection*{A required constraint fails exactly}
Theorem 5.5, equation (5.15), on PDF p. 51/printed p. 979 requires \(H_k^iZ^{\Theta^r}=0\) for \(1\le i\le r\) and \(k\ge-i+2\). In particular \(H_1^2Z^{\Theta^2}=0\).\par
Equation reference: eq:nonzero. For \(Z_0\), all normally ordered terms of [eq:Htwo] (\#eq:Htwo) at \(r=2,k=1\) vanish except \(p_1=p_2=1\). A negative index forces its partner to be at least \(3\), whose derivative annihilates \(Z_0\); zero-index terms vanish because \(J_0=0\). The surviving coefficient is \(-1/2\), and the central term is zero. Consequently \[
 H_1^2Z_0=-\frac{\hbar^2}{2}\partial_{t_1}^2Z_0
 =-\frac{9\hbar^2}{128}(1-t_1)^{-17/8}\ne0.\] At all times zero the value is already \(-9\hbar^2/128\). This is an exact formal coefficient, not a numerical discrepancy. It proves that \(Z_0\) cannot be the descendant potential characterized by the full constraints, despite satisfying every written premise of Proposition 5.6. The retained reviewer-written SymPy calculation confirms [eq:Hzero] (\#eq:Hzero)–[eq:nonzero] (\#eq:nonzero); the finite expansion is supplementary to the analytic identity.\par
\subsection*{Proof failure, core impact, and counterevidence}
At \(r=2\), the range in (5.19) contains only conformal weight \(i=1\). It does not provide \(H_k^2Z=0\) for positive \(k\). The next commutator paragraph on PDF p. 53/printed p. 981 nevertheless sets \([H_k^2,H_0^2]Z=0\) using the preceding constraints. From \(H_0^2Z=\mu Z\) alone this does not follow: the second term involves \(H_0^2(H_k^2Z)\), for which the required annihilation has not been established. The coefficient \(k(1-r)-r+2\) also becomes \(-k\) and vanishes at \(k=0\). Fixing \(\mu=0\) cannot repair the sufficiency claim, because the explicit example already has that value.\par
Proposition 5.18 on PDF pp. 62–63/printed pp. 990–991 explicitly invokes Proposition 5.6 to reduce identification of the BGW function to the string equation. The proof of Theorem 5.20 on PDF p. 63/printed p. 991 then invokes Proposition 5.18. Thus the false endpoint criterion is a stated proof bridge to a principal conclusion, which supports the Major assessment under the assigned mathematical standard.\par
The demonstrated failure has a narrower scope than falsity of that conclusion. No claim is made that \(Z_0\) is a KP or KdV tau function, that the particular BGW function fails a constraint, or that the BGW equality is false. The paper states on PDF pp. 10 and 63 that full \(r=2\) Virasoro constraints were already established in its references [5,33]. Those external proofs were not independently audited here. They are material counterevidence against interpreting this proof defect as a refutation of the final equality.\par
A defensible repair route is to exclude \(r=2\) from the stated universal criterion and treat that endpoint separately using verified full Virasoro constraints and the uniqueness in Theorem 5.5. The remaining \(r\ge3\) proof and the cited replacement dependencies require their own verification; this report does not certify them. Adding only \(Z(0)=1\), the genus expansion, the dimension selection rule, or \(\mu=0\) is insufficient, since the example satisfies all four.\par
\section*{Additional checks and their boundaries}
\subsection*{Selection, rank, and stable domain}
Definitions 2.2 and 2.4, equations (2.10)–(2.12) and (2.21)–(2.24), and Proposition 2.3 occur on PDF pp. 12–14/printed pp. 940–942. The root condition is \[|a|\equiv-(2g-2+n)\pmod r.\] In the operative nonnegative-primary domain, on a smooth curve, \[\deg L=-\frac{2g-2+n+|a|}{r}<0.\] Hence \(h^0(L)=0\), and Riemann–Roch gives \[D=h^1(L)=g-1+\frac{2g-2+n+|a|}{r}
   =\frac{(r+2)(g-1)+n+|a|}{r}.\] The root congruence makes this rank integral. If the congruence fails, the corresponding root component is empty. At \(r=2,a_i=1\), the specialization is \(D=2g-2+n\), in agreement with the dimension calculation used above. This is a smooth-curve rank and domain check. It does not independently reconstruct local freeness and base change over every compactified nodal stratum, and does not claim vanishing for arbitrary negative primary labels.\par
\subsection*{Deformation degree, node matching, and modified unit}
Definition 2.5 and Proposition 2.6, PDF pp. 15–16/printed pp. 943–944, add \(m\) zero-primary markings and then forget them. Before forgetting, the rank is \(D+m/r\); pushforward lowers degree by \(m\), leaving \[D-\frac{r-1}{r}m.\] This explains finiteness and why \(m>0\) cannot alter the degree-\(D\) term. The rendered source on PDF p. 16 reads \(D+m/r\). An earlier extraction lost the denominator and appeared to read \(D+m\); actual image inspection rejects that apparent discrepancy as an extraction artifact.\par
For the nonseparating node with labels \(b,c\), the child rank is \[D+\frac mr-1+\frac{b+c}{r}.\] The inverse Chiodo pairing permits \(b+c=r\) or \(b=c=0\). The first case has the required degree; in the second case the rank is one smaller and the relevant higher Chern component vanishes. The child’s normalization is \(r\) times the parent’s; expressing the raw child class in normalized form cancels the inverse-pairing factor \(r\).\par
Appending a primary label \(r-1\) raises rank by one. The normalization in (2.21) changes by \(-r\), while (2.19) at \(s=-1\) contributes \(-\psi/r\); their product is \(+\psi\), as in the modified-unit equation (2.28). These checked degrees and signs yield no substantive issue. A complete separating-node and stack-theoretic reconstruction was not undertaken.\par
\subsection*{Unmarked exception and finite \(\kappa\)-coefficients}
Theorem B and Corollary C, PDF pp. 5–7/printed pp. 933–935, retain an unmarked top-degree exception. In the homogeneity comparison of Theorem 3.13 and Proposition 3.15, PDF pp. 26–28/printed pp. 954–956, equality of the one-marked expressions gives, for stable \(n=0\) and thus \(g\ge2\), \[[\delta(g-1)-d]\bigl(\Omega_{g,0}^{d}
                   -\widetilde\Omega_{g,0}^{d}\bigr)=0.\] Division is justified only away from \(d=\delta(g-1)\). For \(\delta=3\), this is precisely \(d=3g-3\), preserved in Corollaries 3.24–3.25 on PDF pp. 32–33/printed pp. 960–961. The review does not promote the stated cohomological relations to a Chow-ring or Pixton-relation conclusion.\par
Independently, exact rational expansion of \[-\log\left(\sum_{k\ge0}(-1)^k(2k+1)!!\,u^k\right)
 =\sum_{m\ge1}s_m u^m\] gives \[(s_1,s_2,s_3,s_4,s_5)
 =\left(3,-\frac{21}{2},69,-\frac{2529}{4},
                        \frac{36243}{5}\right).\] Exponentiating \(\sum s_m\kappa_m u^m\) agrees with all displayed weighted-degree components zero through five on PDF pp. 6–7. For example, the first nonconstant components are \[3\kappa_1,\qquad
 \frac32(3\kappa_1^2-7\kappa_2),\qquad
 \frac32(3\kappa_1^3-21\kappa_1\kappa_2+46\kappa_3).\] All six components and the exact calculation are retained. This finite algebraic check does not prove the geometric vanishing relations, Teleman’s theorem, all-degree identities, or the unmarked Euler-characteristic integral.\par
\section*{Coverage, evidence, and disposition}
The actual primary source-image inspection covers PDF pages 1, 5, 6, 7, 9, 10, 12, 13, 14, 15, 16, 26, 27, 28, 32, 33, 48, 49, 50, 51, 52, 53, 62, and 63. The decisive evidence for the single finding is the eight source pages \(48\)–\(53\), \(62\), and \(63\), the exact operator calculation, and the separately attributed native challenge. Their current bytes and hashes are bound in the evidence inventory. Mechanical extraction and a programmatic heading scan are not counted as full-paper reading.\par
The compactified spin geometry, full \(R\)-matrix/vacuum construction, stable-graph expansion, Airy asymptotics, complete topological-recursion correspondence, and all Kac–Schwarz and external-reference arguments were not independently reconstructed. The report establishes the stated counterexample and local positive checks, not an exhaustive certification of the 69-page article. No lack of author code is classified as a mathematical defect.\par
The delivery consists of the source PDF and manifest, this English TeX/PDF report, coverage and findings records, per-issue evidence, exact calculation files, actual build and render logs, page-inspection bindings, and an inventory of current support files. The final build record distinguishes real cache warm-up from the stable `–only-cached` build; successful rendering alone is not recorded as visual inspection. Backend model identity is not inferred from a user-interface label. Provenance preserves the runtime’s declared identity, the controller’s recorded UI observation, and the native helpers’ bounded roles separately.\par
The assessment is Major, for the false \(r=2\) universal implication and its use in a core proof route. This is not a rejection decision or independent acceptance: `scientific\_decision=pending` and `human\_reviews=[]`. The proposed repair is a reviewer suggestion, with no assertion of author adoption.\par
\section*{Recovery and source binding}
Paper: `M25-0090`. Packet SHA-256: `c678d7615841c4f2ec9e0ee7c7066bbfd12fe2fc815448a6b01ad8949469e7d3`. Source PDF: 1,853,309 bytes; SHA-256
`970b06c96dac503883d2be2697c85968b4b16a893e42914aeb87ed3f57c75e45`. Source manifest SHA-256:
`2fafe3b3405c794c4c49a67c45a81e9ecd44174121d27ab56cf4b477337f0ca2`.\par
1 N. K. Chidambaram, E. Garcia-Failde, and A. Giacchetto, Relations on \(\overline{\mathcal M}_{g,n}\) and the negative \(r\)-spin Witten conjecture, Invent. Math. 241 (2025), 929–997. doi:10.1007/s00222-025-01351-y (https://doi.org/10.1007/s00222-025-01351-y). Reviewed source: preserved publisher Version of Record.\par
\end{document}
